%%% FIELD legal-fiber-replacement
\(\mathcal F(\Sigma)\) is the set
\[
\left\{(A,D):A\Sigma+\Sigma A^{\top}=D,\quad D=\operatorname{diag}(d,1-d),\quad 0<d<1,\quad |A_{12}|\ge\eta,\quad |A_{21}|\ge\eta,\quad \lVert A\rVert_F\le M\right\}.
\]
Equivalently, for a raw generator-sign pair \((B,\widetilde D)\) satisfying \(B\Sigma+\Sigma B^{\top}=-\widetilde D\), put \(\lambda=\operatorname{tr}(\widetilde D)>0\), \(A=-B/\lambda\), and \(D=\widetilde D/\lambda\). Its normalized representative belongs to \(\mathcal F(\Sigma)\) exactly when \(\widetilde D\) is strictly positive diagonal and
\[
\frac{|B_{12}|}{\operatorname{tr}(\widetilde D)}\ge\eta,\qquad
\frac{|B_{21}|}{\operatorname{tr}(\widetilde D)}\ge\eta,\qquad
\frac{\lVert B\rVert_F}{\operatorname{tr}(\widetilde D)}\le M.
\]
%%% FIELD published-legal-fiber-construction
For the complete directed two-cycle with self-loops
\[
\mathcal E=\{(1,1),(1,2),(2,1),(2,2)\},
\]
define
\[
\mathcal F_{\mathrm{pub}}(\Sigma)
=\left\{(B,\widetilde D):
\widetilde D\ \text{is diagonal and strictly positive},\quad
B\Sigma+\Sigma B^{\top}=-\widetilde D,\quad
B\ \text{is Hurwitz},\quad
\operatorname{supp}(B)=\mathcal E
\right\}.
\]
%%% FIELD published-polarity-construction
\[
\mathcal P_{\mathrm{pub}}(\Sigma)
=\left\{\operatorname{sign}(B_{12}B_{21}):(B,\widetilde D)\in\mathcal F_{\mathrm{pub}}(\Sigma)\right\}
\subseteq\{-1,+1\}.
\]
%%% FIELD published-class-polarity-statement
For every \(\Sigma=\begin{pmatrix}s&t\\ t&u\end{pmatrix}\in\mathbb S_{++}^{2}\) that is \(t\)-faithful for the complete directed two-cycle \(\mathcal E\), meaning \(t\ne0\), the published legal fiber is nonempty and
\[
\mathcal P_{\mathrm{pub}}(\Sigma)=\{-1,+1\}.
\]
More precisely, after trace normalization, every fixed \(d\in(0,1)\) admits full-support witnesses of both product polarities. Thus removal of the analyst-imposed constants \(\eta\) and \(M\) makes same-sign witnesses reappear for every \(t\ne0\), and the published-class product-polarity set is never a singleton.
%%% FIELD published-class-polarity-proof
Fix such a covariance. Positive definiteness gives \(s>0\), \(u>0\), and \(\Delta=su-t^2>0\), while \(t\)-faithfulness gives \(t\ne0\). Fix any \(d\in(0,1)\) and let \(A(d,q;\Sigma)\) be the affine-fiber matrix. By the affine-fiber lemma,
\[
A\Sigma+\Sigma A^{\top}=D(d),
\qquad D(d)=\operatorname{diag}(d,1-d)\succ0,
\]
and \(-A\) is Hurwitz. Consequently \(B=-A\) and \(\widetilde D=D(d)\) obey
\[
B\Sigma+\Sigma B^{\top}=-\widetilde D,
\]
so only structural minimality and the desired product sign remain to be arranged.

The four entries of \(A\) are
\[
A_{11}=\frac{du/2+qt}{\Delta},\quad
A_{12}=\frac{-dt/2-qs}{\Delta},\quad
A_{21}=\frac{-(1-d)t/2+qu}{\Delta},\quad
A_{22}=\frac{(1-d)s/2-qt}{\Delta}.
\]
The off-diagonal zero values of \(q\) are
\[
q_b=-\frac{dt}{2s},
\qquad
q_c=\frac{(1-d)t}{2u},
\]
and the diagonal zero values are
\[
q_{11}=-\frac{du}{2t},
\qquad
q_{22}=\frac{(1-d)s}{2t}.
\]
All four values are finite. If \(t>0\), then \(q_b<0<q_c\); if \(t<0\), then \(q_c<0<q_b\). Hence the open interval
\[
I_d=\bigl(\min\{q_b,q_c\},\max\{q_b,q_c\}\bigr)
\]
is nonempty. Directly from the two affine formulas, for \(q\in I_d\) the entries \(A_{12}\) and \(A_{21}\) have the same nonzero sign, whereas for
\[
q<\min\{q_b,q_c\}
\quad\text{or}\quad
q>\max\{q_b,q_c\}
\]
they have opposite nonzero signs. Removing the two points \(q_{11}\) and \(q_{22}\) from the nonempty open interval leaves a point, and removing them from either exterior ray also leaves a point. At either selected point all four entries of \(A\), and therefore all four entries of \(B=-A\), are nonzero. Thus \(\operatorname{supp}(B)=\mathcal E\).

The interval choice supplies an element of \(\mathcal F_{\mathrm{pub}}(\Sigma)\) with \(B_{12}B_{21}>0\), and an exterior-ray choice supplies one with \(B_{12}B_{21}<0\). Structural minimality excludes a zero product, so \(\mathcal P_{\mathrm{pub}}(\Sigma)\subseteq\{-1,+1\}\); the two witnesses prove the reverse inclusion. No magnitude lower bound or drift-norm upper bound was used, which proves the final assertion.
%%% FIELD affine-fiber-proof
The trace-one coordinate \(D(d)=\operatorname{diag}(d,1-d)\) is the representative supplied by the time-scale-gauge lemma. For a solution of the Lyapunov equation, set
\[
R=A\Sigma-D/2.
\]
By the Lyapunov-stationarity assumption,
\[
R+R^{\top}=A\Sigma+\Sigma A^{\top}-D=0.
\]
Every real skew-symmetric \(2\times2\) matrix is uniquely \(-Q(q)\) for one \(q\in\mathbb R\). Since \(\Sigma\) is invertible, this gives the unique representation
\[
A=(D/2-Q(q))\Sigma^{-1}.
\]
Using
\[
\Sigma^{-1}=\frac1\Delta\begin{pmatrix}u&-t\\-t&s\end{pmatrix}
\quad\text{and}\quad
D=\operatorname{diag}(d,1-d),
\]
matrix multiplication yields
\[
A(d,q;\Sigma)=\frac1\Delta
\begin{pmatrix}
du/2+qt&-dt/2-qs\\
-(1-d)t/2+qu&(1-d)s/2-qt
\end{pmatrix}.
\]
Conversely, this formula gives \(A\Sigma=D/2-Q(q)\), and hence
\[
A\Sigma+\Sigma A^{\top}=D/2-Q(q)+D/2+Q(q)=D.
\]

It remains to derive stability. By the strict-diffusion assumption, \(0<d<1\), so \(D\succ0\). Define
\[
C=\Sigma^{-1/2}A\Sigma^{1/2}.
\]
The matrix \(C\) is similar to \(A\), and multiplying the Lyapunov identity on both sides by \(\Sigma^{-1/2}\) gives
\[
C+C^{\top}=\Sigma^{-1/2}D\Sigma^{-1/2}\succ0.
\]
If \(Cv=\mu v\) for a nonzero complex vector \(v\), then
\[
2\operatorname{Re}(\mu)\lVert v\rVert_2^2
=v^*(C+C^{\top})v>0.
\]
Thus every eigenvalue of \(C\), and hence of \(A\), has strictly positive real part. Every eigenvalue of \(-A\) therefore has strictly negative real part, so \(-A\) is Hurwitz.
%%% FIELD sector-characterization-statement-replace
For each \(\sigma\), let \(z_{\sigma}\) be the unique minimizer of \(f\) on nonempty \(H_{\sigma}\), with \(v_{\sigma}=f(z_{\sigma})\). The sector contains a legal model if and only if \(H_{\sigma}\ne\varnothing\), \(J_{\sigma}\) holds, and
\[
v_{\sigma}<M^2\quad\text{or}\quad\bigl(v_{\sigma}=M^2\ \text{and}\ 0<(z_{\sigma})_d<1\bigr).
\]
Moreover, for \(\omega\in\{-1,+1\}\), \(H_{(\omega,\omega)}\ne\varnothing\) if and only if \(-\omega t\max(s,u)\ge2\eta\Delta(s+u)\), and \(J_{(\omega,\omega)}\) holds if and only if the inequality is strict or \(s=u\) at equality; both opposite-sign sectors always satisfy \(J_{\sigma}\). Hence \(+1\in\mathcal P(\Sigma)\) exactly when one same-sign sector passes the certificate, and \(-1\in\mathcal P(\Sigma)\) exactly when one opposite-sign sector passes it.
%%% FIELD sector-characterization-proof
The certificate-completeness theorem proves that, whenever \(H_{\sigma}\ne\varnothing\), the retained equality-active family contains the unique minimizer \(z_{\sigma}\) of the strictly convex quadratic \(f\) over \(H_{\sigma}\). It also proves the exact strict-diffusion and norm-attainment criterion
\[
H_{\sigma}\ne\varnothing,\qquad
J_{\sigma},\qquad
v_{\sigma}<M^2
\quad\text{or}\quad
\left(v_{\sigma}=M^2\ \text{and}\ 0<(z_{\sigma})_d<1\right). \tag{1}
\]
The assumptions on the two edge margins and the drift norm are precisely the inequalities encoded in that program, while the legal-fiber definition identifies strict diffusion with \(0<d<1\). Thus (1) proves the first assertion. It remains to derive the advertised scalar rules for \(J_{\sigma}\) and then to aggregate the four sectors by product sign.

Write \(b=A_{12}\) and \(c=A_{21}\). The affine-fiber lemma gives
\[
b=\frac{-dt/2-qs}{\Delta},
\qquad
c=\frac{-(1-d)t/2+qu}{\Delta}.
\]
Consequently the circulation coordinate cancels from the weighted sum:
\[
ub+sc
=-\frac{t\{s+(u-s)d\}}{2\Delta}. \tag{2}
\]
The four signed edge-margin inequalities are, respectively,
\[
\begin{array}{lll}
b\ge\eta &\Longleftrightarrow& q\le(-\eta\Delta-dt/2)/s,\\
b\le-\eta&\Longleftrightarrow& q\ge(\eta\Delta-dt/2)/s,\\
c\ge\eta &\Longleftrightarrow& q\ge(\eta\Delta+(1-d)t/2)/u,\\
c\le-\eta&\Longleftrightarrow& q\le(-\eta\Delta+(1-d)t/2)/u.
\end{array} \tag{3}
\]
Every division in (3) is legitimate because positive definiteness of \(\Sigma\) gives \(s>0\), \(u>0\), and \(\Delta>0\).

Fix \(d\in[0,1]\). For the same-sign sector \((\omega,\omega)\), comparison of the relevant lower and upper endpoints in (3) shows that a feasible \(q\) exists if and only if
\[
-\omega t\{s+(u-s)d\}\ge2\eta\Delta(s+u). \tag{4}
\]
Indeed, for \(\omega=+1\), (4) is obtained by requiring the lower bound for \(c\ge\eta\) not to exceed the upper bound for \(b\ge\eta\); for \(\omega=-1\), it is obtained by requiring the lower bound for \(b\le-\eta\) not to exceed the upper bound for \(c\le-\eta\). Thus no implication is lost in passing from (3) to (4).

The bracket in (4) is the affine function \(s(1-d)+ud\), which is positive on \([0,1]\) and whose maximum is \(\max(s,u)\). Since the right-hand side of (4) is strictly positive, feasibility also forces \(-\omega t>0\). Maximizing the left-hand side over the closed interval therefore gives
\[
H_{(\omega,\omega)}\ne\varnothing
\quad\Longleftrightarrow\quad
-\omega t\max(s,u)\ge2\eta\Delta(s+u). \tag{5}
\]
If the inequality in (5) is strict, continuity of the affine function in (4) gives a satisfying \(d\in(0,1)\). If equality holds and \(s=u\), the bracket is constant, so every \(d\in(0,1)\) satisfies (4) at equality. If equality holds and \(s\ne u\), the bracket reaches its maximum at exactly one of \(d=0\) and \(d=1\), so no interior \(d\) is feasible. This proves the complete stated rule for \(J_{(\omega,\omega)}\), including its open-versus-closed equality branch.

For the opposite-sign sector \((+1,-1)\), (3) imposes two upper bounds on \(q\). For any fixed \(d\in(0,1)\), a real \(q\) smaller than both bounds satisfies the two margins. For \((-1,+1)\), (3) imposes two lower bounds, and a \(q\) larger than both satisfies them. Hence both opposite-sign sectors always satisfy \(J_{\sigma}\).

Finally, the edge-margin assumptions exclude zero off-diagonal entries. A sector with equal edge signs has \(\operatorname{sign}(A_{12}A_{21})=+1\), and a sector with opposite edge signs has product sign \(-1\). By the polarity-set definition, taking the union over the legal sectors and applying (1) proves the two final equivalences.
%%% FIELD joint-only-statement-change
For \(\kappa=1/2\), \(\eta=1/10\), and \(M=1\), put
\[
\Sigma_{\circ}=\begin{pmatrix}1&1/100\\1/100&1\end{pmatrix}.
\]
There exists \(\varepsilon_0>0\) such that the open Frobenius ball
\[
U_I=\left\{\Sigma\in\mathbb S_{++}^{2}:\lVert\Sigma-\Sigma_{\circ}\rVert_F<\varepsilon_0\right\}
\]
is \(t\)-faithful, is contained in \(\mathcal K\), and satisfies \(\mathcal P(\Sigma)=\{-1\}\) for every \(\Sigma\in U_I\), while each of \(A_{12}\) and \(A_{21}\) takes both signs across \(\mathcal F(\Sigma)\). The two sign-reversing witnesses may be chosen with the fixed coordinates \(d=1/2\) and \(q=\pm1/8\), and their drift matrices have full support throughout \(U_I\). This singleton assertion is for the bounded-margin fiber \(\mathcal F\), not for the unrestricted published fiber \(\mathcal F_{\mathrm{pub}}\).
%%% FIELD joint-only-statement-replace
For \(\kappa=1/2\), \(\eta=1/10\), and \(M=1\), put
\[
\Sigma_{\circ}=\begin{pmatrix}1&1/100\\1/100&1\end{pmatrix}.
\]
There exists \(\varepsilon_0>0\) such that the open Frobenius ball
\[
U_I=\left\{\Sigma\in\mathbb S_{++}^{2}:\lVert\Sigma-\Sigma_{\circ}\rVert_F<\varepsilon_0\right\}
\]
is \(t\)-faithful, is contained in \(\mathcal K\), and satisfies \(\mathcal P(\Sigma)=\{-1\}\) for every \(\Sigma\in U_I\), while each of \(A_{12}\) and \(A_{21}\) takes both signs across \(\mathcal F(\Sigma)\). The two sign-reversing witnesses may be chosen with the fixed coordinates \(d=1/2\) and \(q=\pm1/8\), and their drift matrices have full support throughout \(U_I\). This singleton assertion is for the bounded-margin fiber \(\mathcal F\), not for the unrestricted published fiber \(\mathcal F_{\mathrm{pub}}\).
%%% FIELD joint-only-proof
At the displayed center write \(r=1/100\), so \(\Delta=1-r^2=9999/10000\). Its eigenvalues are \(1-r=99/100\) and \(1+r=101/100\); hence
\[
\frac12 I_2\prec\Sigma_{\circ}\prec2I_2,
\]
and \(t=r\ne0\), so the center is \(t\)-faithful.

Fix \(d=1/2\). On the equal-variance slice, the affine-fiber formula becomes
\[
A(d,q;\Sigma_{\circ})
=\frac1\Delta
\begin{pmatrix}
1/4+qr&-r/4-q\\
-r/4+q&1/4-qr
\end{pmatrix}. \tag{6}
\]
For \(q=1/8\), the off-diagonal magnitudes are
\[
|A_{12}|=\frac{1275}{9999}>\frac1{10},
\qquad
|A_{21}|=\frac{1225}{9999}>\frac1{10}, \tag{7}
\]
with \(A_{12}<0<A_{21}\). For \(q=-1/8\), the two magnitudes and signs are interchanged, so \(A_{12}>0>A_{21}\). In either case the diagonal entries are the two positive numbers
\[
\frac{4975}{19998}
\quad\text{and}\quad
\frac{5025}{19998}. \tag{8}
\]
Thus both matrices have full support. Each of their four entry magnitudes is smaller than \(1/3\), and therefore
\[
\lVert A\rVert_F^2<4/9<1. \tag{9}
\]
Equations (7)--(9), strict diffusion, and the affine-fiber lemma show that both choices are legal negative-product witnesses, and across them each directed coefficient takes both signs.

It remains to exclude positive product. The same-sign feasibility condition in the sector-characterization theorem is necessary only if
\[
|t|\max(s,u)\ge2\eta\Delta(s+u). \tag{10}
\]
At \(\Sigma_{\circ}\), the left side of (10) is \(1/100\), whereas the right side is
\[
2\left(\frac1{10}\right)\left(\frac{9999}{10000}\right)(2)
=\frac{9999}{25000}>\frac1{100}. \tag{11}
\]
Thus both same-sign sectors are empty at the center.

All inequalities in the covariance bounds, (7), (9), and the reverse of (10) are strict, and every entry in (8) is nonzero. The affine formula and all quantities in these inequalities are continuous on \(\mathbb S_{++}^{2}\). Hence there is \(\varepsilon_0>0\) such that throughout the displayed Frobenius ball: \(t\ne0\); the covariance bounds remain strict; the two fixed \((d,q)\) choices retain their edge signs, strict margins, full support, and norm below \(M=1\); and the reverse strict inequality to (10) persists. Shrinking \(\varepsilon_0\) if necessary also keeps the ball inside \(\mathbb S_{++}^{2}\). The fixed witnesses prove \(U_I\subseteq\mathcal K\) and \(-1\in\mathcal P(\Sigma)\), while the sector-characterization theorem excludes \(+1\). Therefore \(\mathcal P(\Sigma)=\{-1\}\) throughout \(U_I\), with the two claimed individual-edge sign ambiguities. This is a bounded-margin-refinement result; it does not contradict the unrestricted published-class theorem because the latter imposes neither \(\eta\) nor \(M\).
